\section{Cadrage}

Le secteur du développement web et mobile, en particulier dans un contexte de forte concurrence
entre agences et sociétés de services numériques, fait face à un défi récurrent : la difficulté à connaître
précisément le positionnement tarifaire et la réputation des entreprises concurrentes. Cette problématique
revêt une importance stratégique, car une méconnaissance du marché peut conduire à un mauvais
positionnement commercial, à des offres mal chiffrées, voire à une perte de compétitivité face à des
concurrents mieux informés.

Les défis spécifiques rencontrés dans ce domaine incluent :
\begin{itemize}
    \item \textbf{Dispersion et hétérogénéité de l'information} : les prix et les avis clients des concurrents
    sont disséminés sur de multiples supports (sites vitrines, plateformes d'avis, réseaux sociaux), sans
    format ni structure commune, ce qui rend leur collecte manuelle longue et peu fiable.
    \item \textbf{Évolution constante du marché} : les tarifs pratiqués et la réputation des concurrents évoluent
    en continu, rendant obsolète toute veille ponctuelle réalisée à un instant donné.
    \item \textbf{Absence d'outillage adapté} : les équipes commerciales ne disposent d'aucun outil centralisé
    leur permettant d'interroger rapidement ces informations au moment de chiffrer une offre ou de
    préparer une prospection.
\end{itemize}

Ce projet vise à développer un système capable d'automatiser la collecte de ces données, de les
structurer et de les rendre accessibles à travers un agent conversationnel intelligent, en s'appuyant sur des
techniques d'intelligence artificielle générative et de traitement automatique du langage. L'objectif est de
substituer une démarche manuelle et ponctuelle par une veille continue, structurée et directement
exploitable par des utilisateurs non techniques.

\subsection{Besoins exprimés et opportunités identifiées}

ICHTUS IT a exprimé le besoin de disposer d'une vision consolidée et actualisée du marché sur lequel elle
évolue, sans mobiliser de ressources humaines supplémentaires pour cette veille. Concrètement, l'entreprise
souhaite que ses commerciaux puissent, en une question posée en langage naturel, obtenir à la fois le
positionnement tarifaire d'un ou plusieurs concurrents et une synthèse qualitative de leur réputation
client. Cette opportunité s'inscrit dans la volonté plus large de l'entreprise d'outiller ses équipes
commerciales avec des solutions basées sur la donnée plutôt que sur l'intuition.

\subsection{Données disponibles}

Les données mobilisées dans ce projet proviennent exclusivement de sources publiquement accessibles sur
le web : pages tarifaires publiées par les entreprises concurrentes, et avis clients publiés sur des plateformes
publiques. Ces données sont hétérogènes par nature (texte libre, tableaux de prix, notes chiffrées) et
nécessitent un travail de structuration avant exploitation. Le volume de données dépend directement du
nombre de concurrents suivis et de la fréquence de rafraîchissement du scraping, ce volume étant appelé à
croître progressivement au fil du stage à mesure que de nouveaux concurrents sont ajoutés au périmètre de
veille.

\subsection{Contraintes}

Plusieurs contraintes encadrent la réalisation de ce projet :
\begin{itemize}
    \item \textbf{Contraintes légales} : seules des données publiquement accessibles sont collectées, dans le
    respect des conditions d'utilisation des sites concernés et des principes du RGPD, notamment en évitant
    toute collecte de données à caractère personnel identifiable.
    \item \textbf{Contraintes techniques liées à la structure des sites cibles} : la majorité des sites de
    concurrents ciblés sont construits avec des frameworks JavaScript modernes (rendu côté client), ce qui
    rend leur contenu inaccessible à un scraping classique basé sur des bibliothèques comme BeautifulSoup,
    limitées à l'analyse du HTML statique. Cette contrainte impose le recours à un outil de scraping capable
    d'interpréter le rendu dynamique des pages, en l'occurrence ScrapeGraphAI, qui combine navigation
    automatisée et extraction pilotée par un modèle de langage.
    \item \textbf{Contraintes de coût} : le recours à une solution de scraping basée sur des appels à une API
    (ScrapeGraphAI) engendre un coût par requête, ce qui impose d'optimiser la fréquence de rafraîchissement
    des données et de limiter le nombre d'appels redondants, plutôt que de scraper en continu l'ensemble des
    sources suivies.
    \item \textbf{Contraintes métier} : le système doit rester utilisable par des profils non techniques ; toute
    réponse fournie par l'agent conversationnel doit donc être claire, synthétique et directement actionnable
    par un commercial.
    \item \textbf{Contraintes temporelles} : le développement, l'expérimentation et l'évaluation du système
    doivent être réalisés dans le cadre du stage, ce qui impose de prioriser un périmètre fonctionnel réaliste.
\end{itemize}

\subsection{Périmètre géographique et utilisateurs}

Le périmètre retenu est la \textbf{métropole lilloise et ses environs} (Lille, communes limitrophes,
Hauts-de-France proches), sur des prestations de \textbf{développement web et mobile}. Les utilisateurs
cibles sont les commerciaux et l'encadrement commercial d'ICHTUS IT. On distingue les concurrents
\textit{implantés ou suivis} dans ce périmètre (présents en base) des acteurs simplement cités hors zone,
pour lesquels le système doit s'abstenir ou signaler une couverture insuffisante.

\subsection{Objectifs mesurables et questions techniques}

Les objectifs du projet sont formulés de manière vérifiable :
\begin{itemize}
    \item \textbf{O1 --- Collecte et structuration.} Disposer d'une base de concurrents avec fiches prix
    exploitables (cible observée : environ 69 fiches après filtrage qualité) et d'analyses de sentiment
    pour une part significative de ces fiches.
    \item \textbf{O2 --- Interrogation en langage naturel.} Permettre à un commercial d'obtenir, via une
    question en français, une réponse fondée uniquement sur les fiches retrouvées (architecture RAG).
    \item \textbf{O3 --- Robustesse des gardes-fous.} Rejeter ou signaler les questions hors sujet, hors
    zone géographique et hors couverture métier.
    \item \textbf{O4 --- Évaluation.} Séparer un jeu de développement et un jeu de test indépendant ;
    mesurer le comportement de retrieval et compléter par une grille humaine de qualité de réponse.
\end{itemize}

Questions techniques / hypothèses de travail :
\begin{itemize}
    \item \textbf{Q1.} Un RAG sur fiches structurées (prix + avis) suffit-il à produire des réponses
    commercialement utiles, sans fine-tuning d'un modèle de langue ?
    \item \textbf{Q2.} Un seuil de similarité et un filtre géographique structurel réduisent-ils les
    réponses hors périmètre par rapport à un prompt seul ?
    \item \textbf{Q3.} Une analyse de sentiment zero-shot par LLM est-elle suffisamment fiable, sur un
    échantillon annoté, pour enrichir les fiches sans modèle supervisé dédié ?
\end{itemize}

\subsection{Définition des indicateurs de performance (KPI)}

Afin d'évaluer objectivement la réussite du projet, les indicateurs suivants sont retenus :
\begin{itemize}
    \item \textbf{Couverture de la veille} : nombre de concurrents et de sources suivis de manière automatisée.
    \item \textbf{Fraîcheur des données} : délai moyen entre la mise à jour d'une information chez un
    concurrent et sa disponibilité dans le système (à mesurer ou à déclarer non mesuré si hors délai du stage).
    \item \textbf{Comportement de retrieval du RAG} : taux de questions pertinentes avec résultat, taux de
    rejet hors sujet, Hit Rate@$k$ / MRR sur vérité terrain nominative (diagnostics de recherche).
    \item \textbf{Qualité des réponses générées} : grille humaine (sources, montants/devises, fidélité,
    complétude, abstention, utilité commerciale).
    \item \textbf{Qualité de l'analyse de sentiment} : comparaison à une annotation manuelle de référence
    sur un échantillon d'avis (macro-F1 / matrice de confusion).
    \item \textbf{Gain de temps commercial} : réduction estimée du temps de recherche d'information par
    rapport à une veille manuelle (avec le référent métier si disponible).
\end{itemize}

\subsection{Cahier des charges}

\subsubsection{Exigences fonctionnelles}

\begin{itemize}
    \item \textbf{EF1.} Collecter et mettre à jour les données tarifaires et d'avis de concurrents web/mobile
    du périmètre depuis des sources publiques (Clutch, La Fabrique du Net).
    \item \textbf{EF2.} Stocker les données dans une base relationnelle (Supabase/PostgreSQL) avec historique
    des changements de prix.
    \item \textbf{EF3.} Produire une analyse de sentiment structurée (global, score, points forts/faibles)
    à partir des avis.
    \item \textbf{EF4.} Indexer les fiches en embeddings et répondre à des questions en français via RAG,
    en citant les sources utilisées.
    \item \textbf{EF5.} Appliquer des gardes-fous : hors sujet, hors zone, hors couverture, seuil de
    similarité.
    \item \textbf{EF6.} Exposer le pipeline de façon testable en ligne de commande (et, si le calendrier le
    permet, via une API / interface).
\end{itemize}

\subsubsection{Exigences non fonctionnelles}

\begin{itemize}
    \item \textbf{ENF1 --- Utilisabilité.} Réponses en français, synthétiques, actionnables par un
    commercial non technique.
    \item \textbf{ENF2 --- Fidélité.} Interdiction d'inventer des tarifs ou avis absents des fiches ;
    signaler les informations manquantes.
    \item \textbf{ENF3 --- Coût.} Limiter les appels API (scraping, embeddings, LLM) ; privilégier Gemini
    en développement et un modèle plus capable en démonstration si besoin.
    \item \textbf{ENF4 --- Reproductibilité.} Documenter modèles, paramètres, version du corpus, scripts
    d'évaluation et séparation développement / test.
    \item \textbf{ENF5 --- Conformité.} Données publiques uniquement ; pas de données personnelles
    identifiables ; respect des conditions d'usage des sources.
\end{itemize}

\subsubsection{Hors périmètre (Éval~1 / stage)}

Sauf réalisation explicite et évaluée : génération automatisée de leads commerciaux, pricing
prescriptif (recommandation du prix ICHTUS), couverture nationale hors métropole lilloise, fine-tuning
d'un modèle de langue from scratch.

\subsubsection{Critères d'acceptation}

Un critère est \textbf{satisfait} uniquement s'il est démontrable (script, capture, tableau de résultats
ou annotation).

\begin{table}[H]
\centering
\begingroup
\setstretch{1.15}
\small
\begin{tabularx}{\textwidth}{|c|X|X|}
\hline
\textbf{ID} & \textbf{Critère} & \textbf{Preuve attendue} \\
\hline
CA1 & Au moins une base de fiches prix indexées (embeddings non nuls)
& Comptage SQL / script ; $n$ fiches documenté \\
\hline
CA2 & Une question nominative ramène la fiche attendue dans le TOP-$k$
& Hit Rate sur jeu de test / vérité terrain \\
\hline
CA3 & Une question hors sujet ne ramène pas de fiche concurrent
& Scénarios hors sujet du protocole d'évaluation \\
\hline
CA4 & Une question hors zone (ex.\ Tokyo) est bloquée ou signalée
& Filtre géographique + scénario de test \\
\hline
CA5 & Les montants cités dans la réponse sont présents dans les sources
& Grille humaine (exactitude montants/devises) \\
\hline
CA6 & Jeu de développement distinct du jeu de test ; config figée avant le test
& Fichiers \texttt{dev\_questions} / \texttt{test\_questions} + config figée \\
\hline
CA7 & Échantillon d'avis annoté pour évaluer le sentiment (même réduit)
& Règles d'annotation + scores (ex.\ macro-F1) \\
\hline
CA8 & Pipeline exécutable par un tiers autorisé
& README : installation, variables, commandes de test \\
\hline
\end{tabularx}
\endgroup
\caption{Critères d'acceptation du système (Évaluation~1 et suivi)}
\label{tab:criteres_acceptation}
\end{table}

\subsection{Importance des solutions apportées}

\subsubsection{Amélioration de la prise de décision commerciale}

L'automatisation de la veille concurrentielle permet de transformer des données dispersées et non
structurées en informations directement exploitables. En combinant recherche vectorielle et analyse de
sentiment, le système proposé offre une vision à la fois quantitative (positionnement tarifaire) et
qualitative (perception client) du marché, permettant aux commerciaux d'argumenter leurs offres sur des
bases factuelles plutôt que sur des impressions.

L'impact de cette solution dépasse le cadre strictement opérationnel : en réduisant le temps consacré à la
recherche manuelle d'informations, elle libère du temps commercial pouvant être réinvesti dans la relation
client. Elle contribue également à une meilleure réactivité de l'entreprise face aux évolutions du marché.

\subsubsection{Objectifs et impacts attendus}

Les principaux objectifs de ce projet consistent à développer un système fonctionnel de veille
concurrentielle automatisée, adapté aux spécificités du secteur du développement web et mobile. Ce
développement s'accompagne d'une exigence de fiabilité, tant dans la collecte des données que dans la
pertinence des réponses fournies par l'agent conversationnel.

L'évaluation rigoureuse du système, à travers les indicateurs définis précédemment, constitue un axe
essentiel du projet, permettant de mesurer objectivement son apport par rapport à une démarche de veille
manuelle. Enfin, ce projet ambitionne de constituer un socle réutilisable, l'architecture développée pouvant
être étendue à d'autres secteurs d'activité ou à d'autres types de données concurrentielles au sein de
l'entreprise.